\label{sec:introduction}

\subsection{June 25, 2013}

On June 25, 2013, the Supreme Court decided \textit{Shelby County v.\
Holder}.\footnote{\textit{Shelby County v.\ Holder}, 570 U.S.\ 529 (2013).}
By a 5--4 vote, the Court invalidated the coverage formula in Section~4(b)
of the Voting Rights Act of 1965---the provision that designated which
states and counties were subject to Section~5 preclearance---on the ground
that it was based on data from 1964--1972 and no longer reflected current
conditions sufficient to justify its departure from the equal sovereignty
of the states.

Section~5 itself was not struck. The majority opinion by Chief Justice
Roberts acknowledged that Congress could enact a new coverage formula
based on current data, and it left Section~5's operative requirements
in place. But without a valid coverage formula, no jurisdiction is
covered. Section~5 became a dead letter within hours of the decision.

The practical consequences were swift. Texas, which had been one of
the nine fully covered states, announced the same day that it would
implement a strict voter ID law that had previously been rejected
under preclearance. North Carolina enacted an omnibus voting restrictions
bill within months. Other former covered states moved to implement
redistricting plans that had been blocked or conditioned under
the preclearance process.

\subsection{What the Post-\textit{Shelby} Landscape Looks Like}

Thirteen years after \textit{Shelby County}, the VRA enforcement landscape
for redistricting has stabilized around Section~2. The \textit{Thornburg
v.\ Gingles}\footnote{\textit{Thornburg v.\ Gingles}, 478 U.S.\ 30 (1986).}
three-part test---geographic concentration, political cohesion, and
majority bloc voting---remains the operative standard for redistricting
challenges. Section~2 cases proceed ex post, after maps are enacted,
through federal court litigation that can take years. The structural
asymmetry between Section~5 (ex ante, administrative, faster) and Section~2
(ex post, judicial, slower) is the defining feature of the post-\textit{Shelby}
enforcement environment.

\subsection{Algorithmic Redistricting in the Post-\textit{Shelby} World}

This paper argues that algorithmic redistricting fills part of the gap
that Section~5's demise created. Section~5 required covered jurisdictions
to prove, before implementation, that a proposed change did not retrogress
minority voting strength. This ex ante proof requirement created a paper
trail and a neutral benchmark: the existing plan provided the non-retrogression
baseline. Without Section~5, no such proof is required before enactment,
and no neutral baseline exists unless one is constructed.

The \textsc{bisect} pipeline produces a neutral baseline by construction.
Because it operates without partisan data---it uses Census demographics,
geographic adjacency, and population counts, but not election results or
voter registration data---the \textsc{bisect} plan reflects what a
principled, non-partisan redistricting process would produce. The gap
between a \textsc{bisect} plan's BVAP distribution and an enacted plan's
BVAP distribution measures the portion of minority vote dilution that
is attributable to mapmaker choice rather than geography. This is precisely
the comparison that Section~2 litigation requires---and precisely the
comparison that Section~5 would have caught ex ante.

\subsection{Overview}

Section~\ref{sec:background} describes the pre-\textit{Shelby} Section~5
framework and the \textit{Shelby County} holding.
Section~\ref{sec:methodology} traces Section~2's evolution as the primary
post-\textit{Shelby} enforcement mechanism, including the effects of
\textit{Abbott v.\ Perez} (2018) and \textit{Brnovich v.\ Democratic
National Committee} (2021).
Section~\ref{sec:results} analyzes \textit{Allen v.\ Milligan} (2023)
and state VRA analogs.
Section~\ref{sec:legal} develops the algorithmic redistricting as
transparency substitute argument.
Section~\ref{sec:conclusion} concludes.
